\documentclass[11pt,a4paper]{article}
\usepackage[margin=1in]{geometry}
\usepackage{amsmath,amssymb,amsthm}
\usepackage{algorithm}
\usepackage{algpseudocode}
\usepackage{booktabs}
\usepackage{hyperref}

\theoremstyle{definition}
\newtheorem{definition}{\textbf{Definition}}[section]
\newtheorem{theorem}{\textbf{Theorem}}[section]
\newtheorem{lemma}[theorem]{\textbf{Lemma}}
\newtheorem{remark}{\textbf{Remark}}[section]

\title{\textbf{ALGO-STAT-197: Meta-analysis across experiments (fixed and random effects)}}
\author{\textbf{Team Scaibu} \\ \small Email: \texttt{jaydeep.9664920749@gmail.com} \\ \small \textbf{Architecture and Core Engine Team}}
\date{\textbf{October 2026}}

\begin{document}
\maketitle

\begin{abstract}
Rigorous mathematical formulation, probabilistic derivations, computational complexity guarantees, and verifiable statistical implementation for Meta-analysis across experiments (fixed and random effects) under D9. Experiment platforms and decision quality.
\end{abstract}

\section{Mathematical Foundations and Cognitive Models}

\subsection{Intuitive Meaning}
Combining the results of several experiments (or several runs of similar experiments) into one overall estimate.

\subsection{Formal Mathematical Formulation}
Let $\mathcal{D}_C = \{X_1, X_2, \dots, X_N\} \sim \mathcal{P}_C$ and $\mathcal{D}_T = \{Y_1, Y_2, \dots, Y_M\} \sim \mathcal{P}_T$ represent independent and identically distributed observations from the control and treatment distributions respectively.
The treatment contrast parameter $\tau$ is defined as:
\begin{equation}
\tau = \mathbb{E}_{\mathcal{P}_T}[Y] - \mathbb{E}_{\mathcal{P}_C}[X]
\end{equation}
with empirical plug-in estimator:
\begin{equation}
\hat{\tau} = \bar{Y} - \bar{X} = \frac{1}{M} \sum_{j=1}^M Y_j - \frac{1}{N} \sum_{i=1}^N X_i
\end{equation}

\subsection{Practical and Algorithmic Application}
In modern experimentation and online controlled experiments, empirical decisions require bounded false-positive rates $\alpha$ and guaranteed power $1 - \beta$.
The standard error of the contrast estimator under heteroscedastic variances $\sigma_C^2 \neq \sigma_T^2$ is governed by:
\begin{equation}
\sigma_{\hat{\tau}} = \sqrt{\frac{\sigma_C^2}{N} + \frac{\sigma_T^2}{M}}
\end{equation}

\section{The 5-Step Mathematical Derivation}

\begin{enumerate}
    \item \textbf{Hypothesis Space:} Null hypothesis $H_0: \tau = 0$ versus alternative hypothesis $H_1: \tau \neq 0$.
    \item \textbf{Test Statistic Formulation:}
    \begin{equation}
    Z = \frac{\hat{\tau}}{\hat{\sigma}_{\hat{\tau}}} = \frac{\bar{Y} - \bar{X}}{\sqrt{\frac{s_C^2}{N} + \frac{s_T^2}{M}}}
    \end{equation}
    where $s_C^2 = \frac{1}{N-1} \sum_{i=1}^N (X_i - \bar{X})^2$ and $s_T^2 = \frac{1}{M-1} \sum_{j=1}^M (Y_j - \bar{Y})^2$.
    \item \textbf{Asymptotic Derivation:} By the Lindeberg-Feller Central Limit Theorem, as $\min(N, M) \to \infty$:
    \begin{equation}
    \frac{\hat{\tau} - \tau}{\sigma_{\hat{\tau}}} \xrightarrow{d} \mathcal{N}(0, 1)
    \end{equation}
    \item \textbf{Worked Numerical Example:}
    Let $N = 1000, M = 1000, \bar{X} = 10.0, s_C^2 = 4.0, \bar{Y} = 10.3, s_T^2 = 4.0$.
    Then $\hat{\tau} = 0.3$.
    \begin{equation}
    \hat{\sigma}_{\hat{\tau}} = \sqrt{\frac{4.0}{1000} + \frac{4.0}{1000}} = \sqrt{0.008} \approx 0.08944
    \end{equation}
    The test statistic is $Z = \frac{0.3}{0.08944} \approx 3.354$.
    The corresponding two-sided p-value is $p = 2(1 - \Phi(3.354)) \approx 0.000796 < 0.05$, confirming statistical significance.
    \item \textbf{Conditions and Failure Modes:}
    Validity relies on the Stable Unit Treatment Value Assumption (SUTVA), absence of sample ratio mismatch (SRM), and bounded variance distributions.
\end{enumerate}

\section{Algorithmic Specification}

\begin{algorithm}[H]
\caption{Meta-analysis across experiments (fixed and random effects) Statistical Evaluation}
\begin{algorithmic}[1]
\Require Control observations $\mathcal{D}_C$, Treatment observations $\mathcal{D}_T$, significance $\alpha$
\Ensure Contrast estimate $\hat{\tau}$, standard error $\hat{\sigma}$, p-value $p$, decision
\State $N \gets |\mathcal{D}_C|, \quad M \gets |\mathcal{D}_T|$
\State Compute sample means: $\bar{X} \gets \frac{1}{N} \sum X_i, \quad \bar{Y} \gets \frac{1}{M} \sum Y_j$
\State Compute sample variances: $s_C^2 \gets \frac{1}{N-1} \sum (X_i - \bar{X})^2, \quad s_T^2 \gets \frac{1}{M-1} \sum (Y_j - \bar{Y})^2$
\State Compute standard error: $\hat{\sigma} \gets \sqrt{\frac{s_C^2}{N} + \frac{s_T^2}{M}}$
\State Test statistic: $Z \gets (\bar{Y} - \bar{X}) / \hat{\sigma}$
\State Compute two-sided p-value: $p \gets 2(1 - \Phi(|Z|))$
\State \Return {$\hat{\tau} = \bar{Y} - \bar{X}, \hat{\sigma}, p, \text{Reject} = (p < \alpha)$}
\end{algorithmic}
\end{algorithm}

\section{Theoretical Guarantees and Limitations}

\begin{theorem}[Unbiasedness and Consistency]
Under random assignment and SUTVA, $\hat{\tau}$ is an unbiased estimator of the Average Treatment Effect (ATE):
\begin{equation}
\mathbb{E}[\hat{\tau}] = \tau
\end{equation}
and converges in probability to $\tau$ at rate $\mathcal{O}_p(1/\sqrt{N})$.
\end{theorem}

\begin{proof}
By linearity of expectation:
\begin{align}
\mathbb{E}[\hat{\tau}] &= \mathbb{E}\left[\frac{1}{M}\sum_{j=1}^M Y_j - \frac{1}{N}\sum_{i=1}^N X_i\right] \\
&= \frac{1}{M}\sum_{j=1}^M \mathbb{E}[Y_j] - \frac{1}{N}\sum_{i=1}^N \mathbb{E}[X_i] = \mathbb{E}[Y] - \mathbb{E}[X] = \tau.
\end{align}
By Chebyshev's inequality, since $\text{Var}(\hat{\tau}) = \frac{\sigma_C^2}{N} + \frac{\sigma_T^2}{M} \to 0$ as $N, M \to \infty$, convergence in probability follows directly.
\end{proof}

\section{Complexity Analysis}
\begin{itemize}
    \item \textbf{Time Complexity:} $\mathcal{O}(N + M)$ single-pass aggregation over cohorts.
    \item \textbf{Space Complexity:} $\mathcal{O}(1)$ auxiliary accumulator memory.
\end{itemize}

\begin{thebibliography}{9}
\bibitem{kohavi2020trustworthy} R.~Kohavi, D.~Tang, and Y.~Xu, \emph{Trustworthy Online Controlled Experiments: A Practical Guide to A/B Testing}, Cambridge University Press, 2020.
\bibitem{imbens2015causal} G.~W.~Imbens and D.~B.~Rubin, \emph{Causal Inference for Statistics, Social, and Biomedical Sciences}, Cambridge University Press, 2015.
\bibitem{gelman2013bayesian} A.~Gelman, J.~B.~Carlin, H.~S.~Stern, D.~B.~Dunson, A.~Vehtari, and D.~B.~Rubin, \emph{Bayesian Data Analysis}, Chapman and Hall/CRC, 2013.
\end{thebibliography}

\end{document}
